\chapter{Contrôle et évaluation expérimentale}
\label{ch:control_eval}

\todoF{Chapitre majoritairement à compléter au fil des Sprints 3-5. Structure
proposée ci-dessous, avec ce qu'on a déjà acquis.}

\section{Cinématique directe (\textit{forward kinematics})}
\label{sec:fk}

La cinématique directe calcule la pose de l'effecteur depuis les angles
articulaires :
\begin{equation}
  T_{\text{base}}^{\text{gripper}} = \prod_{i=1}^{n} T_{i-1}^{i}(q_i)
\end{equation}
où $T_{i-1}^{i}(q_i)$ est la transformation d'un repère articulaire au
suivant, paramétrée par l'angle $q_i$.

\paragraph{Implémentation.} Le module
\texttt{src/calibration/forward\_kinematics.py} parse l'URDF
\texttt{configs/so101\_new\_calib.urdf} (récupéré du dépôt officiel
TheRobotStudio/SO-ARM100) avec \texttt{xml.etree.ElementTree} et compose la
chaîne cinématique en NumPy. La chaîne SO-101 comprend 5 articulations
rotoïdes actionnées (\texttt{shoulder\_pan}, \texttt{shoulder\_lift},
\texttt{elbow\_flex}, \texttt{wrist\_flex}, \texttt{wrist\_roll}) plus un
joint fixe vers le repère \texttt{gripper\_frame\_link}.

\paragraph{Décision D4 (FK en NumPy stdlib).} L'alternative envisagée
(\textit{placo} + \textit{pinocchio}, stack utilisée par LeRobot) a été
écartée : trop lourde (bindings C++, met à jour numpy), boîte noire pour la
rédaction. Réimplémentation en NumPy pur : $\sim$80 lignes, auto-testée,
explicable dans le mémoire. L'URDF reste la \textbf{seule source de vérité
géométrique}.

\section{Cinématique inverse (\textit{inverse kinematics})}
\label{sec:ik}

\todoF{À compléter au Sprint 3, brique 3.2. Plan prévu :
solveur IK numérique via \texttt{scipy.optimize.least\_squares}
(Levenberg-Marquardt), avec contraintes de plage articulaire. Voir
\texttt{src/control/ik.py} (à créer).}

\section{Génération de trajectoire}
\label{sec:trajectory}

\todoF{À compléter au Sprint 3, brique 3.3. Plan prévu : interpolation
trapézoïdale dans l'espace articulaire entre les 3 poses (approach, grasp,
retract). Vitesse maximale et accélération maximale paramétrables pour le
SO-101.}

\section{Module de contrôle moteur}
\label{sec:motor_control}

\todoF{À compléter au Sprint 3, brique 3.4. Plan prévu : wrapper autour de
\texttt{FeetechMotorsBus} avec envoi de commandes \texttt{sync\_write
Goal\_Position} à fréquence régulière, gestion du retry en cas d'erreur
bus.}

\section{Pipeline intégré}
\label{sec:pipeline}

\todoF{À compléter au Sprint 3, brique 3.5. Plan prévu : script
\texttt{src/pipeline.py} qui orchestre :\\
\textbf{boucle} : capture caméras $\rightarrow$ détection $\rightarrow$
triangulation 3D $\rightarrow$ grasp planning $\rightarrow$ IK $\rightarrow$
trajectoire $\rightarrow$ exécution.}

\section{Évaluation expérimentale}
\label{sec:eval}

Les quatre métriques imposées par le cahier des charges (objectif 6) sont :

\begin{description}
  \item[\textbf{Taux de réussite de la saisie}] Sur N essais avec un objet
        donné, fraction des essais où le robot saisit et place correctement
        l'objet dans la boîte de dépose.
  \item[\textbf{Nombre de replans}] Nombre de re-planifications déclenchées
        par la boucle fermée de raffinement (détection d'écart significatif entre
        position estimée et position réelle pendant le mouvement).
  \item[\textbf{Collisions évitées}] Nombre de fois où la trajectoire
        contourne un obstacle de la scène.
  \item[\textbf{Précision de la prise}] Erreur de position de l'effecteur au
        moment du \textit{grasp}, mesurée \textit{a posteriori} via la
        cinématique directe sur la pose moteur réelle.
\end{description}

\subsection{Protocole de mesure}

\todoF{À détailler au Sprint 5. Plan : $\sim$20 essais par configuration
expérimentale (objet seul, objet entouré, objet partiellement caché),
$\sim$5 objets cibles, $\sim$3 positions différentes par objet. Total
$\sim$300 essais. Analyse statistique : moyenne $\pm$ écart-type, intervalle
de confiance 95\%.}

\subsection{Comparaison V1 (modulaire) vs V2 (imitation learning)}

\todoF{Sprint 5 : enregistrer $\sim$50 démonstrations téléopérées pour 1
tâche, entraîner SmolVLA via \texttt{lerobot-train}, comparer les 4 métriques
ci-dessus entre les deux approches sur les \textbf{mêmes scènes}. Tableau
attendu :}

\begin{center}
\begin{tabular}{lccc}
\toprule
Métrique & Modulaire (V1) & Modulaire (V2-HF) & Imitation (SmolVLA) \\
\midrule
Taux réussite (\%) & \texttt{??} & \texttt{??} & \texttt{??} \\
Nb replans moyens & \texttt{??} & \texttt{??} & \texttt{??} \\
Collisions évitées & \texttt{??} & \texttt{??} & \texttt{??} \\
Précision prise (mm) & \texttt{??} & \texttt{??} & \texttt{??} \\
\bottomrule
\end{tabular}
\end{center}
